\addcontentsline{toc}{subsubsection}{Solution 64. The Invisible Conjugate}

\begin{dlsummary}[title={\bfseries Solution 64. The Invisible Conjugate},breakable]

\begin{proof}
The only solutions are
\[
\boxed{f(a+b\sqrt2)=|a+b\sqrt2|}
\]
and
\[
\boxed{f(a+b\sqrt2)=|a-b\sqrt2|.}
\]

\paragraph{1. Establish symmetry and the reverse triangle inequality.}

Multiplicativity gives
\[
f(-1)^2=f(1)=1.
\]
Nonnegativity implies \(f(-1)=1\), hence
\[
f(-x)=f(x).
\]
Subadditivity now gives
\[
\boxed{|f(x)-f(y)|\le f(x-y).} \tag{1}
\]

\paragraph{2. Determine the function on the rationals.}

Repeated addition gives
\[
f(n)\le n\qquad(n\in\mathbb Z_{>0}),
\]
while multiplicativity gives \(f(2^k)=2^k\).

For any positive integer \(n\), choose \(2^k>n\). Then
\[
2^k=f(2^k)
\le f(n)+f(2^k-n)
\le f(n)+2^k-n.
\]
Thus \(f(n)\ge n\), proving equality.

Multiplicativity and symmetry consequently yield
\[
\boxed{f(q)=|q|\qquad(q\in\mathbb Q).} \tag{2}
\]
Also,
\[
f(\sqrt2)^2=f(2)=2,
\]
so \(f(\sqrt2)=\sqrt2\).

\paragraph{3. Force one conjugate direction to vanish.}

Choose rational numbers \(r_n\to\sqrt2\), and define
\[
A_n=f(\sqrt2-r_n),\qquad
B_n=f(\sqrt2+r_n).
\]

By (1) and (2),
\[
|A_n-A_m|\le|r_n-r_m|,
\qquad
|B_n-B_m|\le|r_n-r_m|.
\]
Therefore both sequences converge; write their limits as \(A,B\ge0\).

Multiplicativity gives
\[
A_nB_n
=f\bigl((\sqrt2-r_n)(\sqrt2+r_n)\bigr)
=f(2-r_n^2)
=|2-r_n^2|.
\]
Hence
\[
AB=0.
\]
But subadditivity gives
\[
A_n+B_n\ge f(2\sqrt2)=2\sqrt2,
\]
so
\[
A+B\ge2\sqrt2.
\]
Thus \textbf{exactly one} of \(A,B\) is zero.

\paragraph{4. The vanishing direction determines every value.}

If \(A=0\), then, for any rational \(a,b\), equation (1) gives
\[
\begin{aligned}
\left|f(a+b\sqrt2)-|a+br_n|\right|
&=\left|f(a+b\sqrt2)-f(a+br_n)\right|\\
&\le f\bigl(b(\sqrt2-r_n)\bigr)\\
&=|b|A_n.
\end{aligned}
\]
Taking limits,
\[
f(a+b\sqrt2)=|a+b\sqrt2|.
\]

If \(B=0\), compare instead with the rational number \(a-br_n\):
\[
\left|f(a+b\sqrt2)-|a-br_n|\right|
\le |b|B_n.
\]
Taking limits gives
\[
f(a+b\sqrt2)=|a-b\sqrt2|.
\]

This proves completeness.

\paragraph{5. Verify both solutions.}

The first is ordinary absolute value.

For the second, define
\[
\overline{a+b\sqrt2}=a-b\sqrt2.
\]
Direct expansion shows
\[
\overline{x+y}=\bar x+\bar y,
\qquad
\overline{xy}=\bar x\,\bar y.
\]
Therefore \(f(x)=|\bar x|\) satisfies both functional conditions and all normalisations.

\paragraph{A revealing example.}

At \(x=1+\sqrt2\), the two solutions give
\[
f(x)=1+\sqrt2
\quad\text{or}\quad
f(x)=\sqrt2-1.
\]
So an input larger than \(2\) can have a function value smaller than \(1\).
\end{proof}
\end{dlsummary}

\begin{takeaways}
\begin{itemize}
\item \textbf{Identify a familiar subdomain first:} Exact values on \(\mathbb Q\) provide the approximation machinery.
\item \textbf{Conjugate products can force a choice:} Two nonnegative limiting quantities have product zero but cannot both vanish.
\item \textbf{A local choice can determine everything:} Once one direction becomes small, rational comparison fixes all inputs.
\item \textbf{Do not assume ordinary continuity:} One solution is discontinuous relative to the usual real-number distance on \(K\).
\end{itemize}
\end{takeaways}
